\chapter{揭示思维链推理对直接偏好优化的影响：来自Text-to-SQL的教训}

\subsection{背景知识：偏好优化}
偏好优化是一种让模型与人类偏好对齐的技术，它通常需要一个包含“好/坏”回答对比的训练集。最经典的方法是\textbf{基于人类反馈的强化学习（Reinforcement Learning from Human Feedback, RLHF）}，它先训练一个奖励模型来模拟人类偏好，再用强化学习（如PPO）优化语言模型。然而，RLHF流程复杂且不稳定。

\textbf{直接偏好优化（Direct Preference Optimization, DPO）}用一个巧妙的重参数化技巧，把奖励模型直接融进了损失函数，从而绕过了强化学习步骤，让偏好优化能像监督学习一样简单稳定。DPO的目标是最大化正确回答与错误回答之间的对数似然差异，同时确保模型参数不会偏离原始模型太远。核心损失函数如下：
\[
\mathcal{L}_{\mathrm{DPO}}(\pi_{\theta}; \pi_{\mathrm{ref}}) = -\mathbb{E}_{(x,y_w,y_l)\sim\mathcal{D}} \left[ \log \sigma \left( \beta \log \frac{\pi_{\theta}(y_w|x)}{\pi_{\mathrm{ref}}(y_w|x)} - \beta \log \frac{\pi_{\theta}(y_l|x)}{\pi_{\mathrm{ref}}(y_l|x)} \right) \right]
\]
其中 \(\pi_{\theta}\) 是训练中的模型，\(\pi_{\mathrm{ref}}\) 是参考模型（通常是SFT后的模型），\(y_w\) 和 \(y_l\) 分别代表正确（偏好）和错误（被拒绝）的回答，\(\beta\) 是个超参数，控制对偏离参考模型的惩罚强度，\(\sigma\) 是sigmoid函数。

这个公式是DPO（直接偏好优化）的核心损失函数，它的设计目的是让模型学会更倾向于生成“好”的回答、远离“坏”的回答，同时保证模型不会偏离原本的参考模型太远。下面我们逐部分拆解它的含义。本质上是在做对比学习，拉近模型与好回答的距离，推远与坏回答的距离。但与普通对比学习不同，这里的“距离”是用相对于参考模型的概率比来度量的。公式中的\(\beta \log \frac{\pi_{\theta}(y|x)}{\pi_{\mathrm{ref}}(y|x)}\)可以看作模型为输出\(y\)分配的隐式奖励。损失函数就是在最大化好回答与坏回答之间的奖励差。由于参考模型\(\pi_{\mathrm{ref}}\)固定，优化过程中模型如果为了增大差值而过度调整参数，会导致概率比变得极大，但\(\beta\)的存在会抑制这种偏离（因为差值中包含了\(\pi_{\mathrm{ref}}\)的对数概率，模型无法随意增加好回答的概率而不影响其他输出）。

传统RLHF需要先训练一个奖励模型，再用PPO等强化学习算法优化语言模型，流程复杂且不稳定。DPO直接通过这个损失函数，将偏好优化转化为监督学习问题，训练过程简单、高效，且效果往往不逊于RLHF。这个公式通过比较当前模型与参考模型在好坏回答上的概率差异，引导模型在保持基本能力的同时，学会区分优劣输出，从而实现与人类偏好对齐。它是DPO算法的核心，也是实现偏好优化“简单化”的关键。

此后，研究者们又提出了许多DPO的变体，比如SimPO、KTO、IPO等，它们在不同场景下各有千秋。在本研究中，研究人员采用最经典的DPO作为代表算法，因为它效果已经过广泛验证，而且实现简单。值得一提的是，目前只有一项研究（SENSE, Yang et al., 2024c）尝试把DPO用到Text-to-SQL上。不过他们的方法跟研究人员有本质区别：SENSE是用小模型（1B）为大模型生成偏好数据，而研究人员遵循标准DPO流程，直接从SFT模型采样构建偏好数据。更重要的是，SENSE仍然用SQL查询作为训练标签，而研究人员的研究则探索了将CoT解决方案作为训练标签的关键作用。

\subsection{从执行反馈中学习}
在LLM火起来之前，通过执行反馈进行学习已经在代码生成任务中崭露头角。比如用单元测试结果作为奖励信号，通过强化学习优化代码生成模型。类似的工作还有RLTF (Liu et al., 2023)、StepCoder (Dou et al., 2024) 等。在Text-to-SQL领域，执行反馈天然可用——因为SQL查询可以在数据库上执行并与标准结果对比。然而，现有工作要么没有充分利用执行反馈，要么在在线学习方式中面临巨大的计算开销（执行SQL需要大量内存和CPU资源）。本文是首个通过离线构建偏好数据（利用执行反馈标记正确/错误）并结合DPO实现稳定性能提升的工作。

\section{引言}

近年来，Text-to-SQL在自然语言处理和数据库研究领域掀起了一股热潮。随着大型语言模型（LLM）的横空出世，解决Text-to-SQL的路径主要分成了两条。

第一条是\textbf{基于提示的方法}：直接向GPT-4、Gemini这样的闭源巨头提问，让它们生成SQL。理论可行，但实际操作中往往需要设计复杂的提示，甚至动用多个智能体协作，成本高、响应慢，而且把数据送到外部服务还可能引发隐私担忧。

第二条是\textbf{监督微调（SFT）方法}：在公开的Text-to-SQL数据集（比如Spider、Bird）上对开源模型进行微调，让它们学会把问题映射到SQL。微调后的模型可以本地部署，成本低、速度快，数据也安全。可惜的是，SFT的性能被高质量训练数据的稀缺性卡住了脖子——创建这些数据既烧钱又耗时，模型只能模仿训练集里的“标准答案”，一旦遇到没见过的复杂问题就容易露馅了。

为突破监督微调的瓶颈，也许可以借助一种叫“偏好优化”的算法来增强SFT模型的表现，如\textbf{直接偏好优化（Direct Preference Optimization, DPO）}。和SFT只让模型学正确答案不同，DPO会喂给模型一些“偏好对”——比如一个正确回答和一个错误回答，然后训练模型更倾向于生成正确的，同时远离错误的。这招在数学解题和代码生成这类需要复杂推理的任务中屡试不爽。但在Text-to-SQL领域却几乎没人用，这是为什么呢？

\subsection{碰壁}

为了找到答案，研究人员在一个颇具挑战性的Text-to-SQL基准——Bird数据集上进行试验。Bird的每个样本都包含一个问题、一个数据库和对应的标准SQL查询。挑10个参数规模从6.7B到15B的开源LLM，按标准DPO流程一步步来：先做监督微调，让模型学会从问题和数据库信息生成SQL；然后用微调好的模型为每个训练样本采样多个SQL查询，再到数据库上执行，把结果正确的标记为“正确”，错误的标为“错误”，配对成偏好对；最后用DPO损失在这些偏好对上进行一通优化。

实验结果却不尽如人意。DPO压根没能稳定提升性能——10个模型里，有6个的准确率反而下滑。研究人员试遍了各种挽救的招数：调超参数、往DPO损失里加SFT损失、把DPO换成其他偏好优化算法（比如KTO）、甚至学别人的方法用小模型生成偏好数据。可折腾一圈下来，最好的情况也不过是不到1.5\%的微小提升。

\subsection{反思}
撞了南墙之后，研究人员开始反思：为什么DPO在数学和代码任务中屡试不爽，到了Text-to-SQL就水土不服了呢？对比两类任务的数据集，研究人员发现了它们之间的一个关键差异。数学和代码数据集，比如MATH、GSM8K、CodeUltraFeedback，通常不只给最终答案，还会附上详细的解题步骤或推理过程，也就是所谓的“思维链”（Chain-of-Thought, CoT）解决方案。思维链推理之所以能释放DPO的潜力，是因为它能缓解“奖励破解”（模型为了得高分而走捷径），增强模型区分好坏的能力，并让训练过程更好地规模化扩展。这些推理过程就像一座桥梁，把问题和最终答案连接起来，让模型在SFT和DPO训练中能更好地理解背后的逻辑，从而举一反三。

反观Text-to-SQL数据集，像Bird、Spider、WikiSQL，它们都只给最终答案——标准的SQL查询语句，中间怎么一步步推出来的完全没有。模型在SFT阶段只能死记硬背SQL模板，在DPO阶段也只能在最终SQL层面比对比对，根本缺乏对推理过程的监督。这就像只给学生答案而不教他们解题思路。那么如果为Text-to-SQL数据自动补上思维链，应该就会实现DPO在Text-to-SQL任务上稳定且显著的性能提升。

\subsection{验证猜想：给数据补上CoT}
为了验证这个想法，研究人员设计了一个流程来研究CoT究竟如何影响DPO在Text-to-SQL中的表现。研究人员找来一个LLM充当“CoT合成器”，让它为现有的Text-to-SQL数据集自动生成逐步的CoT解决方案。这个合成器会接收数据库信息、问题以及标准SQL查询作为输入，然后输出详细的推理步骤，一步步展示如何将问题转化为SQL。用这种方法，研究人员以极低的人力成本构建了CoT增强的Bird数据集。

接着，研究人员在完全相同的实验设置下，对CoT增强的数据集进行了SFT和DPO训练。结果如图1所示：\textbf{加上CoT之后，所有10个评估模型在DPO阶段都获得了显著且稳定的性能提升！}研究人员还不放心，又把评估范围扩展到Spider、SpiderDK、Spider-Syn等多个Text-to-SQL基准，结果一致，无一例外。

\subsection{刨根问底：CoT到底起了什么作用？}
为了搞清楚CoT为什么这么重要，研究人员深入分析，总结出三个核心发现。第一，引入CoT显著减少了DPO训练中常见的“奖励破解”现象——也就是模型通过生成看似合理但实际错误的输出来骗取高分，让训练变得更加稳定有效。第二，CoT增强了DPO作为隐式奖励模型的判别能力，让它能更准确地判断一个回答是好是坏。第三，CoT提高了DPO的可扩展性，无论增加偏好数据量，还是在推理时加大采样预算，CoT增强的DPO都能持续受益。

\section{方法论}
本节将详细介绍研究人员提出的流程，它包含三个环环相扣的步骤：思维链（CoT）合成、监督微调（SFT）、直接偏好优化（DPO）。图2展示了整个流程的概览。

\subsection{思维链合成}
现有的Text-to-SQL数据集，比如Bird和Spider，都只提供最终答案——标准SQL查询，而没有从问题到SQL的中间推理过程。为了让模型学会像人一样一步步思考，研究人员首先需要为每个训练样本生成CoT解决方案。

为什么非得要CoT呢？因为对于一个复杂的Text-to-SQL问题，人类解决它通常需要经历一系列步骤：先理解问题到底在问什么，再识别出涉及哪些数据库表和列，然后考虑有什么过滤条件，需不需要聚合，最后才能构建出完整的SQL语句。如果模型也能学习这样的中间步骤，它就能更好地理解问题的逻辑，而不是简单地死记硬背SQL模板。更重要的是，在DPO训练中，有了CoT，模型就可以在词元级别获得更细粒度的反馈，而不是只在最终SQL层面判断对错，这能有效缓解奖励破解的问题。

那怎么生成CoT？为了尽量减少人工标注的成本，研究人员请来一个LLM当“CoT合成器”。具体来说，对于训练集中的每个样本（包括问题 \(q\)，数据库信息 \(d\)，标准SQL \(s\)），研究人员用GPT-4o-mini（OpenAI 2024b）生成 \(K\) 个多样化的CoT解决方案。每个CoT解决方案 \(c\) 都是一段文本，它会一步步展示如何从问题出发，结合数据库模式，最终推导出正确的SQL查询。因为最终答案（标准SQL）已经提供给了合成器，所以生成的CoT解决方案既可靠又准确，完全可以放心地用于后续的SFT和DPO训练。需要注意的是，研究人员并不要求所有生成的CoT都完美无瑕，但通过引导模型生成多样化的推理路径，可以为后续训练提供丰富的语义信息。在本工作中，研究人员设定 \(K=16\)，也就是为每个原始样本生成16个CoT解决方案。

\subsection{直接偏好优化（DPO）}
SFT虽然能让模型学会生成CoT，但它本质上还是在模仿训练数据，可能仍然无法区分好的推理路径和坏的推理路径。DPO则通过偏好数据进一步优化模型，让它更倾向于生成正确的回答。

\subsubsection{偏好数据集构建}
研究人员需要构建一个偏好数据集 \(D_P\)，其中每个样本包含一个问题 \(q\)、数据库信息 \(d\)、一个正确的CoT解决方案 \(c^+\) 和一个错误的CoT解决方案 \(c^-\)。构建过程是这样的：
首先，对于训练集中的每个样本，研究人员用参考模型 \(\pi_{\mathrm{SFT}}\) 生成 \(N\) 个不同的CoT解决方案（研究人员取 \(N=16\)）。接着，从每个生成的CoT解决方案里，用正则表达式提取出预测的SQL查询。

然后，把预测SQL和标准SQL分别在对应的数据库上执行，比较它们的执行结果（比如返回的记录是不是完全一致）。如果结果相同，就把这个CoT解决方案标记为“正确”；否则就标记为“错误”。最后，对于每个训练样本，如果它同时有至少一个正确和一个错误的CoT解决方案，研究人员就随机挑选一个正确和一个错误的，组成一个偏好对 \((c^+, c^-)\)，并把它加入偏好数据集 \(D_P\)。
因为并不是每个样本都能同时产生正确和错误的CoT（比如模型可能全对或全错），最终得到的偏好数据集大小通常小于SFT数据集（以Bird为例，大约有1.5k到2.5k对）。

\subsubsection{DPO学习目标}
DPO的目标是让正确回答相对于错误回答的“优势”最大化，同时确保模型不会跑偏太远，离参考模型太远。它隐含地定义了一个奖励模型。对于任意输出 \(y\)，它的隐式奖励可以这样算：
\[
r_{\theta}(x,y) = \beta \log \frac{\pi_{\theta}(y|x)}{\pi_{\mathrm{ref}}(y|x)}
\]
这个奖励可以理解为模型对当前输出相对于参考模型的偏好程度。因为语言模型是自回归的，研究人员还可以把这个奖励分解到每个词元上：
\[
r_{\theta}(x,y) = \sum_{t=1}^{T} \beta \log \frac{\pi_{\theta}(y_t|x,y_{<t})}{\pi_{\mathrm{ref}}(y_t|x,y_{<t})}
\]
这种分解让研究人员能在词元级别分析模型对每个生成步骤的“信心”，对于理解模型行为（比如奖励破解）特别有帮助。举个例子，如果模型在某个词元上给出了异常高的奖励，那可能意味着它找到了一条捷径，值得警惕。在后续的分析中，研究人员会用这个隐式奖励来评估DPO模型的判别能力和训练稳定性。

\section{实验设置}
为了验证研究人员的猜想，并全面评估CoT增强DPO的效果，研究人员精心设计了一系列实验。

\subsection{数据集}
研究人员的实验主要在以下几个Text-to-SQL基准上进行。首先是\textbf{Bird (Li et al., 2024b)}，这是一个颇具挑战性的跨领域数据集，包含9,428个训练样本和1,534个开发样本，涉及37个专业领域的95个大型数据库。Bird的数据库通常拥有大量真实数据行，因此执行结果反馈更为可靠，也更贴近实际应用场景。其次是\textbf{Spider (Yu et al., 2018)}，它是最常用的Text-to-SQL基准，包含7,000个训练样本和1,034个开发样本，涵盖138个领域的200个数据库。

Spider的SQL查询相对简单，数据库行数较少，执行结果反馈可能存在假阳性问题。为了评估模型的泛化能力，研究人员还在Spider的几个变体上进行了测试，包括Spider-DK、Spider-Syn 、Spider-Realistic以及 Dr.Spider。这些数据集通过对Spider开发集进行各种扰动（如同义词替换、数据库结构调整等），模拟了真实世界中的各种挑战。

\subsection{评估指标}
研究人员采用\textbf{执行准确率（Execution Accuracy, EX）}作为主要评估指标。EX衡量的是预测SQL查询与标准SQL查询在数据库上执行后，返回的结果是否完全一致。对于Spider开发集，研究人员还额外报告\textbf{测试套件准确率（Test Suite Accuracy, TS）}。TS通过在多个测试套件数据库实例上执行来减少假阳性，提供更稳健的评估。

\subsection{推理策略}
给定一个训练好的Text-to-SQL模型，研究人员用三种不同的推理策略来评估它的性能。

第一种是\textbf{贪心（Greedy）}，也就是使用温度为0的解码，每次选择概率最高的词元，生成单个确定性的响应。这是最常用的快速评估方式。

第二种是\textbf{Pass@1}，研究人员使用温度为1.0采样一个响应，重复16次，计算平均准确率。这反映了模型在单次采样下的期望性能。

第三种是\textbf{Maj@K}，研究人员使用温度为1.0采样 \(K\) 个响应（通常 \(K=16\)），然后基于每个响应提取的SQL查询在数据库上的执行结果进行多数投票，将得票最多的查询作为最终预测。这种方法通常能提升准确率，因为它聚合了多次采样的共识。

\subsection{实现细节}
在基础模型的选择上，研究人员挑了10个有代表性的开源LLM，参数范围从6.7B到15B，涵盖了不同领域：通用模型有Deepseek-llm-7b-chat、Meta-Llama-3.1-8B-Instruct、Qwen2.5-7B-Instruct、Qwen2.5-14B-Instruct；代码模型有Deepseek-coder-7b-instruct-v1.5、CodeQwen1.5-7B-Chat、Qwen2.5-Coder-7B-Instruct；SQL专用模型有CodeS-7B和CodeS-15B 。这些模型涵盖了当前主流的中小规模开源LLM，让研究人员的结论更具普适性。

所有实验都在配备8张NVIDIA A800 GPU（80GB内存）的服务器上进行，使用CUDA 12.1.1。训练框架采用Llama Factory 并启用了FlashAttention 2.0  和 DeepSpeed ZeRO-3加速训练、减少显存占用。对于超过10B参数的模型，在DPO阶段研究人员启用了CPU卸载以容纳优化器状态。

训练超参数方面，研究人员采用全参数微调，使用bf16混合精度。优化器是AdamW，\(\beta_1=0.9\)，\(\beta_2=0.99\)。学习率调度采用余弦衰减，前5\%的步骤进行线性预热。批次大小统一设为64，上下文窗口为4096个词元。学习率在SFT阶段，对于小于10B的模型设为 \(1\times10^{-5}\)，大于10B的设为 \(7\times10^{-6}\)；在DPO阶段则分别设为 \(1\times10^{-6}\) 和 \(7\times10^{-7}\)。DPO的超参数 \(\beta\) 默认设为0.1（后面研究人员还会专门分析它的敏感性）。训练轮数上，SFT阶段训练4个epoch，研究人员根据开发集上的Maj@16指标选出最佳检查点作为DPO的参考模型；DPO阶段继续训练8个epoch。

数据构建也有一些讲究。输入提示的平均长度为965个词元，里面包含了问题、数据库模式信息（表名、列名、主外键）以及通过模式链接和值匹配检索出的相关表和列值。研究人员沿用了CodeS 的方法：先用一个模式项分类器（基于XLMRoBERTa-XL，3.5B）识别最相关的表和列，然后通过“粗到细”的匹配方法从数据库中检索与问题相关的值，最后组合成提示。输出标签方面，原始数据集的输出是标准SQL查询，平均长度44个词元；而CoT增强后，输出的CoT解决方案平均长度404个词元（词元数由Qwen2.5-7B-Instruct的分词器测量）。CoT合成研究人员使用GPT-4o-mini-2024-07-18，每个样本生成16个CoT解决方案，提示模板见附录B。偏好数据构建则用SFT模型（\(\pi_{\mathrm{SFT}}\)）为每个训练样本生成16个响应，根据执行结果标记正确/错误，然后随机配对形成偏好对。对于Bird，最终偏好数据集大小约为1.5k到2.5k对，具体数量取决于不同模型。

\section{实验结果}
研究人员首先在Bird基准上展示主要发现，然后深入错误分析，最后在Spider及其变体上验证结论。

\subsection{主要结果：Bird开发集}
研究人员对比了两种训练方式：一种是\textbf{Vanilla}，即直接在原始Bird数据上进行SFT和DPO；另一种是\textbf{Syn CoT}，即在CoT增强数据上进行SFT和DPO。对于每种方式，研究人员报告了SFT和DPO阶段在三种推理策略（Greedy, Pass@1, Maj@16）下的执行准确率（EX），以及“Syn CoT + DPO”相对于“Vanilla + SFT”的增益（\(\Delta \mathrm{EX}\)）。

第一个观察是：原始模型在DPO阶段很难获得性能提升。大多数模型在DPO后的准确率变化微乎其微，甚至出现下降。这说明DPO并没能让模型变得更可靠。第二个观察则让人眼前一亮：使用合成CoT的模型在DPO阶段获得了稳定且显著的增益。再看Syn CoT部分，所有模型在DPO后都有了明显进步，说明模型不仅更准，也更可靠了。第三个观察是：合成CoT加DPO展现出了更高的性能上限。所有Syn CoT模型的最终性能都远超对应的Vanilla + SFT基线。

\subsection{错误分析：CoT如何帮助DPO纠正错误？}
为了更深入地理解DPO对模型生成的影响，研究人员以Qwen2.5-7B-Instruct为例，用贪心解码对它在开发集上的预测错误进行了细致的分类。研究人员把错误分成了6大类17小类，每类的详细描述见附录I。表3展示了每个错误类型在DPO后的修正率（Fix \%），也就是DPO后正确而SFT时错误的比例。

研究人员发现，DPO特别擅长纠正那些因忽略细节而引起的错误。修正率超过25\%的错误类型在表中用粗体标出，比如NULL/DISTINCT（用CoT修正了40.0\%）和列顺序（用CoT修正了42.9\%）。这类错误往往是因为模型没有仔细处理数据中的空值、重复值，或者没有按问题要求返回正确的列顺序。而CoT的逐步推理迫使模型在生成SQL之前先考虑这些细节，从而帮助DPO更好地识别和纠正。

CoT还能提供逻辑指导，增强DPO的错误纠正能力。对于涉及多表连接的JOIN错误，Syn CoT DPO的修正率达到了32.1\%，比Vanilla DPO高出了16.5个百分点。这类错误通常需要理解表之间的关系，CoT的逐步逻辑有助于理清连接路径。相比之下，对于幻觉（Hallucination）和日期处理（Date）这类直观错误，CoT的增益就小一些（分别+3.5\%和+7.3\%），因为原始DPO已经能处理得不错了。

\subsection{Spider及其变体的结果}
总体的趋势和Bird一致：Syn CoT + DPO在Spider开发集以及各种鲁棒性基准上都显著优于Vanilla + DPO，证明了这个方法有很好的泛化能力。不过研究人员也注意到，Spider上DPO的增益略小于Bird，这可能是因为Spider的SQL查询相对简单，而且数据库行数少导致执行反馈不够精确（容易出现假阳性）。

接下来，研究人员以Qwen2.5-7B-Instruct为对象，在Bird基准上进行了一系列分析，试图弄清楚CoT对DPO的关键作用。研究人员从三个维度来深入探讨：判别能力、训练稳定性、可扩展性。

\subsection{更好的判别能力}
DPO有个重要特性，就是它隐式学习了一个奖励模型，可以用来评估生成回答的好坏。一个优秀的DPO模型应该能准确地区分正确和错误的回答。为此，研究人员构建了一个评估偏好数据集，其中每个样本包含一个正确和一个错误的CoT解决方案（对于Syn CoT模型）或者对应的SQL（对于Vanilla模型）。然后研究人员计算DPO模型根据隐式奖励选择正确响应的准确率。Syn CoT模型的准确率稳定在较高水平（接近90\%），而Vanilla模型的准确率波动较大，而且整体偏低。这表明CoT推理让DPO模型能够更可靠地判断回答的质量。

\subsection{更好的训练稳定性}
训练稳定性是衡量DPO效果的一个重要指标。研究人员考察了两个现象：奖励破解和超参数敏感性。

先看奖励破解。奖励破解（Reward Hacking）是指模型在训练过程中学会通过生成一些看似合理但实际错误的输出来获取高奖励，从而欺骗奖励模型。在DPO中，研究人员可以通过观察模型在开发集上的真实性能与模型自我评估的隐式奖励之间的差异来检测奖励破解。

对于原始模型，随着训练进行，模型在开发集上的准确率先升后降（在第2个epoch达到峰值后就开始下降），但隐式奖励却一路走高。这说明模型觉得自己的输出越来越好，可实际上性能在下滑，这正是奖励破解的典型表现。而对于Syn CoT模型，准确率先升后趋于饱和，隐式奖励也随之先升后平稳，和真实性能高度一致。这表明CoT有效防止了模型走捷径，让它学习到的奖励与真实性能对齐了。

研究人员进一步分析了模型输出的统计特性。原始模型在DPO后，平均SQL字符长度从1732猛增到1917（增长了26.6\%），同时不可执行SQL的比例从2.0\%飙升到19.6\%（增长了17.6\%）。这说明模型试图通过生成长而复杂的SQL来获取高奖励，但这些SQL往往根本没法执行。相反，Syn CoT模型的SQL长度只微增了2\%，不可执行率反而从7.6\%降到了5.1\%，说明它的输出更可靠。

\subsection{更好的可扩展性}
研究人员还评估了Syn CoT DPO在偏好数据规模（采样预算）和推理时计算（采样数量）方面的可扩展性。图5a显示，对于原始模型，增加偏好数据收集时的采样预算（也就是从每个训练样本生成更多响应来构建偏好对）并不会提升DPO性能；而对于Syn CoT模型，性能却随着预算增加而持续提升。图5b展示了推理时采用Maj@K策略时，增加采样数K对性能的影响。Syn CoT DPO在K较小时（比如K=2）就能达到较高的准确率，而原始DPO需要K=16才能接近这个水平。这表明CoT让模型能够在有限的推理预算下产生更可靠的输出。

\section{前景：集成到现有Text-to-SQL框架中}
研究人员训练的DPO模型可以作为一个更强的基座模型，集成到现有的多智能体Text-to-SQL框架中，进一步提升整体性能。

研究人员选取Qwen2.5-7B-Instruct作为基础模型，把它集成到两个有代表性的框架里。一个是\textbf{DTS-SQL (Pourreza and Rafiei, 2024b)}，它把Text-to-SQL分解成模式链接和SQL生成两个阶段，研究人员替换了它SQL生成阶段的模型。另一个是\textbf{C3-SQL (Dong et al., 2023)}，它包含了清晰提示、模型偏差校准和一致性输出三个模块，研究人员替换了它的核心生成模型。

用Syn CoT + DPO模型替换原框架的基座模型后，两个框架在Spider开发集上的准确率都有了提升，这证明了研究人员训练的模型确实有实用价值。

\section{结论}
回顾这项工作的整个探索历程，研究人员首次证明了通过合成思维链（CoT）解决方案，可以在Text-to-SQL任务上实现直接偏好优化（DPO）的一致且有效的性能改进。通过全面的实验和深入的分析，研究人员揭示了一个关键发现：CoT推理对于释放DPO在Text-to-SQL中的潜力至关重要，因为它能缓解奖励破解、增强模型判别能力、提高训练可扩展性。这些见解不仅解释了为什么先前将DPO直接应用于Text-to-SQL会失败，也为未来构建更强大的开源Text-to-SQL模型指明了方向。研究人员相信，CoT与偏好优化的结合将为复杂推理任务带来更深远的影响。

\section{局限性}
尽管取得了显著的成果，研究人员也很清楚这项研究还存在一些局限性。首先是执行时间问题。Bird数据集中的数据库通常很大，执行SQL查询非常耗时，尤其是在偏好数据构建阶段需要执行大量采样查询。研究人员虽然通过多进程并行执行加快了速度，但多进程间的资源争用可能导致部分查询超时，使一些本应正确的SQL被误判为错误，引入假阴性，这可能会影响DPO的训练质量。其次是反馈不准确性的问题。Spider数据集则相反，因为数据库行数少，执行结果匹配容易产生假阳性（比如两个不同的SQL返回相同的结果），导致偏好信号不可靠。未来可以考虑结合测试套件（如TS）来构建更准确的反馈信号。另外，CoT合成质量也是一个值得探讨的方向。研究人员主要用GPT-4o-mini作为合成器，虽然成本低而且效果显著，但更高质量的CoT（比如人工标注或者用更强模型合成）可能会带来进一步提升。同时，研究人员也没有深入探讨CoT的风格（比如语言、详细程度）对DPO的影响，这可以作为未来研究的一个方向。最后是模型规模的问题。本研究主要聚焦于6.7B到15B参数的中小规模模型，对于更大模型（比如30B以上）是否同样适用还有待验证。不过研究人员相信，方法论本身应该是可以扩展的。

